\begin{proposition}[A ramified conductor with vanishing nearby homology]\hypertarget{tag-02WQ}{}\label{night:P4.19}\label{stmt:P4.19}
\StatementRecord{night:P4.19}{proved}{02WQ}{\SeriesRef{night:P4.16}, \SeriesRef{night:P4.18}}

For the reduced hypersurface germ $y(x^2-zy^2)=0$, the normalization
conductor has multiplicity two on each component and its reduced
map to the singular locus is ramified. The affine nearby fibre is
$\mathbb C^*\times\mathbb C$. Every compact spherical class
in an open subset of this affine fibre has self-intersection zero.
\end{proposition}
\begin{proof}

\begin{enumerate}[label={$\langle1\rangle$\arabic*.},leftmargin=*]
\item One normalization component is the plane $y=0$; the other is
the Whitney normalization $(u,v)\mapsto(uv,v,u^2)$.
The singular locus is the smooth $z$-axis. The plane conductor
is $x^2=0$, obtained by restricting the other component's
equation. Write $R_B=\mathbb C\{v,uv,u^2\}$ and
$N_B=\mathbb C\{u,v\}$. The module $N_B$ is generated over
$R_B$ by $1,u$. A convergent series lies in $R_B$ precisely when
its odd powers of $u$ have coefficients divisible by $v$.
Thus $h,uh\in R_B$ is equivalent to $h\in vN_B$, proving
that the intrinsic conductor is $vN_B$.
The union ring is the fibre product of the component rings over
$\mathbb C\{x,z\}/(x^2)$. On the Whitney component, a
conductor element $h$ must satisfy $hN_B\subset vR_B$; the
intrinsic conductor calculation gives $h\in v^2N_B$.
On the plane, the corresponding kernel is $(x^2)$.
These are the two claimed conductor ideals.
The reduced conductor cover is the identity sheet from the
plane together with $u\mapsto u^2$ from the Whitney component.
The latter is ramified.

\item Put $F=\{y(x^2-zy^2)=1\}$. Here $y\ne0$, and the inverse maps
\[
 (x,y,z)\longmapsto(y,x),\qquad
 (y,x)\longmapsto\left(x,y,\frac{yx^2-1}{y^3}\right)
\]
identify this fibre with $\mathbb C^*\times\mathbb C$.
The retraction $(y,x)\mapsto(y/|y|,0)$ is a homotopy equivalence
with $S^1$, so $\pi_2$ and $H_2$ vanish. For an open inclusion
$j:U\hookrightarrow F$, compact-support duality identifies the
homology map $j_*$ with extension by zero $j_!$ on dual classes.
For the compact-support dual $\alpha$ of a compact spherical class,
\[
 \int_U\alpha^2=\int_F(j_!\alpha)^2=0.
\]
This proves the asserted zero self-intersection in $U$.

\end{enumerate}
\end{proof}